% =====================================================================
% preamble.tex —— 第七—八章合集的统一导言区
% 由各分册的导言合并而成：只保留一份，避免同一宏定义多次重定义。
% =====================================================================

%=====================================================================
% theme-violet.tex —— 讲义精讲 · 紫花主题
% 用法：\documentclass[11pt]{ctexart}
%       \input{assets/theme-violet}
%       \begin{document} ... \end{document}
% 只用 ctexart 自带能力 + xcolor/enumitem/fancyhdr/pifont/amsmath，
% 不依赖 tikz/tcolorbox，保证 MiKTeX 可直接编译。
%=====================================================================

\usepackage[a4paper,left=2.0cm,right=2.0cm,top=2.1cm,bottom=2.3cm,%
            headsep=12pt,footskip=22pt]{geometry}
\usepackage{amsmath,amssymb}
\usepackage{needspace}
\usepackage{xcolor}
\usepackage{array}
\usepackage{enumitem}
\usepackage{pifont}
\usepackage{fancyhdr}
\usepackage[hidelinks]{hyperref}

%---------------------------------------------------------------- 配色
\definecolor{violet}{RGB}{107,78,143}      % 主色：紫
\definecolor{violetdeep}{RGB}{63,44,92}    % 深紫：标题/强调
\definecolor{petal}{RGB}{183,159,208}      % 花瓣：装饰线、花饰
\definecolor{bandviolet}{RGB}{247,243,251} % 浅紫底：档案/提示框
\definecolor{bandcream}{RGB}{252,249,242}  % 米白底：易错/注意框
\definecolor{rulepurple}{RGB}{217,207,230} % 分隔线
\definecolor{inkgray}{RGB}{92,96,110}      % 次要文字

%---------------------------------------------------------------- 字体
% 有 Noto CJK 用 Noto，没有就退回 ctex 默认（Windows 下为宋体/黑体）
\IfFontExistsTF{Noto Serif CJK SC}{\setCJKmainfont{Noto Serif CJK SC}}{}
\IfFontExistsTF{Noto Sans CJK SC}{\setCJKsansfont{Noto Sans CJK SC}}{}

%---------------------------------------------------------------- 版式
\setlist{nosep,leftmargin=1.9em}
\linespread{1.28}
\setlength{\parskip}{3.5pt}
\setlength{\parindent}{0pt}
\setlength{\tabcolsep}{5pt}
\renewcommand{\arraystretch}{1.28}

%---------------------------------------------------------------- 花饰
% \ding{95} = 五瓣小花，\ding{96} = 星形花；都来自 pifont，必定存在
\newcommand{\orn}{\textcolor{petal}{\ding{95}}}          % 行内小标
\newcommand{\ornstar}{\textcolor{petal}{\ding{96}}}      % 星形花
\newcommand{\flowerrule}{%                                % 居中花饰分隔线
  \par\vspace{7pt}%
  \noindent\makebox[\linewidth]{\color{petal}%
    \rule{0.20\linewidth}{0.6pt}\hspace{8pt}{\small\ding{95}}\hspace{4pt}%
    \raisebox{1.1pt}{\tiny\ding{95}}\hspace{4pt}{\small\ding{95}}\hspace{8pt}%
    \rule{0.20\linewidth}{0.6pt}}%
  \par\vspace{7pt}}

%---------------------------------------------------------------- 标题与节头
\newcommand{\doctitle}[3]{%   #1 主标题  #2 副标题  #3 元信息（可用 \\ 分行）
  \begin{center}
    {\color{petal}\small \orn\enspace #2\enspace \ornstar}\\[7pt]
    {\LARGE\bfseries\color{violetdeep} #1}\\[10pt]
    {\color{petal}\rule{0.55\linewidth}{1pt}}\\[4pt]
    {\small\color{inkgray} #3}
  \end{center}\vspace{2pt}}

\newcommand{\hdr}[2]{%   #1 章节号（如 §0）  #2 标题
  \par\vspace{12pt}%
  \noindent\colorbox{violetdeep}{\color{white}\bfseries\small\;#1\;}\hspace{7pt}%
  {\color{violetdeep}\large\bfseries #2}\par\nobreak
  \vspace{3pt}{\color{petal}\hrule height 1.2pt}\vspace{6pt}}

\newcommand{\subhdr}[1]{%
  \par\vspace{6pt}%
  \noindent{\color{violet}\bfseries #1}\par\vspace{1pt}}

%---------------------------------------------------------------- 提示框
\newcommand{\callout}[1]{% 浅紫底：讲义档案、读法说明
  \par\medskip\noindent
  \fcolorbox{rulepurple}{bandviolet}{\parbox{\dimexpr\linewidth-2\fboxsep-2\fboxrule}{\vspace{2pt}#1\vspace{2pt}}}%
  \par\medskip}

\newcommand{\warnbox}[1]{% 米白底：易错、边界条件
  \par\medskip\noindent
  \fcolorbox{rulepurple}{bandcream}{\parbox{\dimexpr\linewidth-2\fboxsep-2\fboxrule}{\vspace{2pt}#1\vspace{2pt}}}%
  \par\medskip}

\newcommand{\gapitem}[2]{% #1 缺口类型（紫） #2 内容
  \par\vspace{2pt}\noindent
  \hangindent=1.4em\hangafter=1
  {\color{petal}\small\ding{96}}\enspace{\bfseries\color{violetdeep}#1}\enspace #2\par}

% 题干块（\qstem）：每道例题／作业的「详解」开头必须先复述题干与出处。
% 起因：2026-10-25 第五章讲义把七道作业的题面集中放在一页、详解另起一页，
% 读者翻到「第 4 题」时看不到条件，等于一道没有题干的题目。
% 规则：详解块内自带题干，读者不翻页就能读懂题在问什么。
\newcommand{\qstem}[1]{%
  \par\medskip\noindent
  \fcolorbox{violet}{bandviolet}{%
    \parbox{\dimexpr\linewidth-2\fboxsep-2\fboxrule}{%
      \vspace{2pt}%
      {\small\bfseries\color{violet}题干}\hspace{0.8em}#1%
      \vspace{2pt}}}%
  \par\medskip}

% v2: 保护题干与开头，长题可跨页，续页补上下文。
\newcommand{\qbreak}{\needspace{6\baselineskip}}
% 独立练习：ID、来源、题干、考点、答案标签（通常 ans:ID）。
\newcommand{\practice}[5]{%
  \qbreak\subhdr{练习 #1}\qstem{（#2）#3}%
  \textbf{考点}：#4\par
  详解见附录第~\pageref{#5}~页。\par\medskip}
% 附录开头：ID、来源、条件摘要；自动产生 ans:ID 标签。
\newcommand{\answerstart}[3]{%
  \qbreak\subhdr{#1\quad 详解}\label{ans:#1}\qstem{（#2）#3}}
\newcommand{\answercontinue}[2]{%
  \needspace{5\baselineskip}\subhdr{#1\quad 续解}%
  {\small\color{inkgray}本页接续：#2}\par\medskip}
% 延迟提供默认值，让旧讲义可在导言区保留自己的 \newcommand{\step}。
\AtBeginDocument{\providecommand{\step}[1]{\needspace{3\baselineskip}\par\medskip\textbf{#1}\par}}

% 前提旁注：只放在「需要理解的那一步」后面，给陈述 + 本题代入，不重证。
\newcommand{\aside}[1]{%
  \par\vspace{1pt}\noindent
  \hangindent=1.6em\hangafter=1
  {\small\color{inkgray}%
    \hspace*{0.6em}{\color{petal}$\triangleright$}\hspace{2pt}#1}%
  \par\vspace{2pt}}

%---------------------------------------------------------------- 表格
% 本机 MiKTeX 的 colortbl 与 array 版本不匹配（\rowcolor 会报错），
% 故表格统一用「上粗紫线 + 表头 + 中细线」的三线式，不依赖 colortbl。
\newcommand{\tabtop}{\par\vspace{3pt}{\color{violet}\hrule height 1.2pt}\vspace{1pt}}
\newcommand{\tabmid}{\par\vspace{1pt}{\color{petal}\hrule height 0.5pt}\vspace{1pt}}
\newcommand{\tabbot}{\par\vspace{3pt}{\color{rulepurple}\hrule height 0.5pt}\vspace{8pt}}
\newcommand{\thead}[1]{{\bfseries\color{violetdeep}#1}}

%---------------------------------------------------------------- 页眉页脚
\pagestyle{fancy}
\fancyhf{}
\renewcommand{\headrulewidth}{0.4pt}
\renewcommand{\headrule}{\hbox to\headwidth{\color{rulepurple}\leaders\hrule height \headrulewidth\hfill}}
\fancyhead[L]{\small\color{inkgray}\songti \doctag}
\fancyhead[R]{\small\color{inkgray}\songti \coursetag}
\fancyfoot[C]{\small\color{inkgray}\orn\enspace{\color{violet}\thepage}\enspace\orn}

% 每篇讲义在导言区设这两个值
\newcommand{\doctag}{}
\newcommand{\coursetag}{}
\newcommand{\settags}[2]{\renewcommand{\doctag}{#1}\renewcommand{\coursetag}{#2}}

%---------------------------------------------------------------- 数学记号
\newcommand{\vect}[1]{\boldsymbol{#1}}
\newcommand{\dd}{\,\mathrm{d}}
\newcommand{\ddx}{\frac{\mathrm{d}}{\mathrm{d}x}}
\addtolength{\jot}{2pt}



\usepackage{tikz}
\usetikzlibrary{arrows.meta,calc,positioning,decorations.pathmorphing}

\settags{数学物理方法（II）· 第七—八章}{精讲讲义 · 2026-11-17}
\hypersetup{pdftitle={数学物理方法（II）第七—八章精讲讲义},
            pdfsubject={第七章 数学物理方程与定解条件；第八章 分离变量法（直角、极坐标、柱坐标）},
            pdfkeywords={数学物理方法, 分离变量, 定解条件, 贝塞尔函数}}

% ---------------------------------------------------------------- 编号与标题
% 三级标题（如 7.1.1）：介于 \hdr（节）与 \subhdr（四级）之间
\newcommand{\hdrsub}[2]{%
  \par\vspace{10pt}\needspace{4\baselineskip}%
  \noindent
  \colorbox{petal}{\color{white}\bfseries\small\;#1\;}\hspace{7pt}%
  {\color{violetdeep}\bfseries\large #2}\par\nobreak
  \vspace{2pt}{\color{rulepurple}\hrule height 0.6pt}\vspace{5pt}}

% 章标题：另起一页，居中，带花饰
\newcommand{\chapterhead}[2]{%
  \clearpage
  \par\vspace*{1.4cm}
  \begin{center}
  {\color{petal}\small \orn\enspace\ornstar\enspace\orn}\\[10pt]
  {\color{violetdeep}\bfseries\Huge #2}\\[8pt]
  {\color{inkgray}\large #1}\\[10pt]
  {\color{petal}\rule{0.46\linewidth}{1.2pt}}
  \end{center}
  \vspace{0.9cm}}

% 组标题：8.1 里的「一维 / 二维」分组，不带编号
\newcommand{\grp}[1]{%
  \par\vspace{11pt}\needspace{3\baselineskip}%
  \noindent{\color{violetdeep}\bfseries\large\orn\enspace #1}\par\nobreak
  \vspace{3pt}{\color{petal}\hrule height 1pt}\vspace{5pt}}

% 标题不落页底
\let\origsubhdr\subhdr
\renewcommand{\subhdr}[1]{\par\needspace{3\baselineskip}\origsubhdr{#1}}
\let\orighdr\hdr
\renewcommand{\hdr}[2]{\par\needspace{6\baselineskip}\orighdr{#1}{#2}}

% ---------------------------------------------------------------- 表格与列
\newcommand{\tabmidrule}{\noalign{\vspace{1pt}{\color{petal}\hrule height 0.5pt}\vspace{1pt}}}
\renewcommand{\thead}[1]{\leavevmode{\bfseries\color{violetdeep}#1}}
\newcolumntype{L}[1]{>{\raggedright\arraybackslash}p{#1}}

% ---------------------------------------------------------------- 图
\newenvironment{figbox}{\par\begin{center}\begin{minipage}{0.96\linewidth}\centering}{\end{minipage}\end{center}\par}
\newcommand{\figcap}[1]{\par\vspace{4pt}{\small\color{inkgray}#1}\par}

% ---------------------------------------------------------------- 定理与证明
\newcommand{\thm}[2]{\par\medskip\noindent{\color{violetdeep}\bfseries #1}\hspace{0.5em}{\bfseries #2}\par\nobreak\vspace{2pt}}
\newcommand{\pfline}{\par\smallskip\noindent\textbf{证明。}\enspace}
\newcommand{\qedsym}{\hfill$\blacksquare$}

% ---------------------------------------------------------------- 单位与圈码
\newcommand{\unit}[1]{\ifmmode\mathrm{#1}\else$\mathrm{#1}$\fi}
\xeCJKDeclareCharClass{CJK}{"2460 -> "24FF}

% ---------------------------------------------------------------- 详解块
\renewcommand{\qbreak}{\needspace{0.30\textheight}}

% ---------------------------------------------------------------- 题解用的宏
% 来自 2026-11-09 的《p19 作业四题中英答案》，附录 B 沿用它的版式
\newcommand{\ok}{\textcolor{violet}{\ensuremath{\checkmark}}}
\newcommand{\en}[1]{\textit{#1}}
\newcommand{\tagline}[1]{{\small\color{inkgray}#1}}
\newcommand{\explain}[1]{%
  \par\medskip\noindent
  \fcolorbox{rulepurple}{bandcream}{\parbox{\dimexpr\linewidth-2\fboxsep-2\fboxrule}{%
    \vspace{2pt}{\color{violetdeep}\bfseries 讲解}\quad #1\vspace{2pt}}}%
  \par\medskip}
\newcommand{\answerbox}[2]{%
  \par\medskip\noindent
  \fcolorbox{violet}{bandviolet}{\parbox{\dimexpr\linewidth-2\fboxsep-2\fboxrule}{%
    \vspace{2pt}{\small\bfseries\color{violet}答案}\hspace{0.8em}#1\par\vspace{3pt}%
    {\small\bfseries\color{violet}Answer}\hspace{0.8em}\en{#2}\vspace{2pt}}}%
  \par\medskip}
\newcommand{\rodfig}[2]{%
  \begin{center}\small
  \begin{tabular}{c@{}c@{}c}
  #1 & {\setlength{\fboxsep}{0pt}\fcolorbox{petal}{petal!45}{\rule{0pt}{0.38cm}\hspace{2.6cm}}} & #2\\[1pt]
   & {\color{inkgray}\makebox[2.6cm][s]{$x=0$\hfill$x=l$}} &
  \end{tabular}
  \end{center}}

% ---------------------------------------------------------------- 页面细节
\setlength{\headheight}{22pt}
\setlength{\abovedisplayskip}{7pt plus 2pt minus 2pt}
\setlength{\belowdisplayskip}{7pt plus 2pt minus 2pt}
\setlength{\abovedisplayshortskip}{4pt plus 1pt}
\setlength{\belowdisplayshortskip}{5pt plus 1pt minus 1pt}
